\%============================================================
\chapter{Bai Hoc: Vuot Qua Nghich Canh --- Su Gap Kho La Giao Vien Tot Nhat (Overcoming Adversity --- Hardship Is the Greatest Teacher)}
\begin{center}
  \textit{\color{truyengold}``Da quy khong mai giua kho khan, va con nguoi khong nen duoc ma khong co thu thach.''}\\[4pt]
  \textit{(The gem cannot be polished without friction, nor man perfected without trials.)}\\[4pt]
  \small\color{truyenblue}--- Co Kim Dong Tay ---
\end{center}
\ngancach

\section{Helen Keller --- Dem Toi Khong The Dap Tat Anh Sang}
\begin{truyen}{Helen Keller --- Dem Toi Khong The Dap Tat Anh Sang}{Helen Keller --- Hoa Ky 1887}
\chuhoa{H}{}H elen Keller bi mat ca thi giac lan thinh giac vao nam 19 thang tuoi vi benh tat.\\
\textit{(Helen Keller lost both her sight and hearing at 19 months of age due to illness.)}

Trong nam thu bay cua cuoc doi minh, ba song trong mot the gioi hoan toan toi va im lang --- khong ngon ngu, khong hinh anh.\\
\textit{(For the first six years of her life she lived in a world of complete darkness and silence --- no language, no images.)}

Ro roi Anne Sullivan den --- va trong mot buoi chieu tai gieng nuoc, Helen dat tay duoi voi nuoc chay, con Anne viet 'n-u-o-c' len long ban tay kia.\\
\textit{(Then Anne Sullivan arrived --- and one afternoon at the water pump, Helen held one hand under flowing water as Anne spelled 'w-a-t-e-r' into her other palm.)}

Helen nho lai: 'Dieu ki dieu da xay ra --- mot cai gi do sat luong sinh dan sang trong toi --- toi biet ray van vat deu co ten.'\\
\textit{(Helen later wrote: 'A miracle happened --- something alive and illuminated woke within me --- I knew that everything had a name.')}

Ba tot nghiep dai hoc, viet 12 quyen sach, du lich 39 quoc gia va tro thanh bieu tuong cua tinh than con nguoi.\\
\textit{(She graduated from college, wrote 12 books, visited 39 countries, and became a symbol of the human spirit.)}

\end{truyen}
\begin{baihoc}
  Bot mot giac quan, con nguoi thuong phat trien nhung nang luc ma nguoi binh thuong khong bao gio can phai kham pha.\\
  \textit{(Deprived of one sense, people often develop capacities that the ordinary person never needs to discover.)}
\end{baihoc}

\section{Beethoven Diec Ma Van Sang Tac}
\begin{truyen}{Beethoven Diec Ma Van Sang Tac}{Ludwig van Beethoven --- Duc 1801}
\chuhoa{K}{}K hoan 27 tuoi, Beethoven bat dau nhan ra minh dang dan bi diec --- mot thoa ma chua chac the nao trai qua.\\
\textit{(Around age 27, Beethoven began realizing he was gradually going deaf --- a fate no composer should have to face.)}

Ong viet trong thu: 'Mo cua re ti toi den gap roi di --- roi moi nguoi di ca --- toi bi bo lai mot minh.'\\
\textit{(He wrote in a letter: 'People open and close doors near me --- they all went away --- I was left alone.')}

Ong can nhac viec tu tu --- nhung quyet dinh tiep tuc vi con qua nhieu am nhac van con trong tam tri ong.\\
\textit{(He considered suicide --- but decided to continue because there was still too much music left in his mind.)}

Bản giao huong so 9 --- tac pham vi dai nhat cua ong --- duoc sang tac khi ong diec hoan toan.\\
\textit{(Symphony No. 9 --- his greatest work --- was composed when he was completely deaf.)}

Khi bieu dien buoi ra mat, ong ngoi tren san khau nhung khong nghe gi ca --- ca khan gia vo tay, ong chi biet khi nguoi ban quay mat ong lai.\\
\textit{(At its premiere he sat on stage but heard nothing --- the audience applauded, he only knew when someone turned him around to face them.)}

\end{truyen}
\begin{baihoc}
  Nhung gi ban mat khong bao gio bang nhung gi ban van con --- Beethoven mat thinh giac nhung khong mat am nhac.\\
  \textit{(What you lose is never equal to what you still have --- Beethoven lost his hearing but never lost his music.)}
\end{baihoc}

\section{J.K. Rowling --- 12 Lan Bi Tu Choi}
\begin{truyen}{J.K. Rowling --- 12 Lan Bi Tu Choi}{J.K. Rowling --- Anh 1995}
\chuhoa{N}{}N am 1995, J.K. Rowling la nguoi me don than, song nho tro cap xa hoi, nuoi con mot minh sau cuoc hon nhan vo.\\
\textit{(In 1995, J.K. Rowling was a single mother on welfare, raising her daughter alone after a failed marriage.)}

Trong hoan canh do, ba viet xong ban thao cuon sach dau tien ve mot cau be phap thuat ten Harry Potter.\\
\textit{(In those circumstances, she finished the first manuscript about a young wizard named Harry Potter.)}

Ba gui ban thao den 12 nha xuat ban --- ca 12 tu choi; mot trong so ho con gop y la khong nen viet cho tre em nua.\\
\textit{(She sent the manuscript to 12 publishers --- all 12 rejected it; one even advised her to stop writing for children.)}

Cuoi cung, nha xuat ban thu 13 --- Bloomsbury --- chap nhan, tuy nguoi chu bien khong cu ki vong cao vao tu be.\\
\textit{(Finally, the 13th publisher --- Bloomsbury --- accepted it, though the editor did not have high commercial hopes.)}

Sach ban duoc hon 500 trieu ban tren toan the gioi, dich ra 80 ngon ngu, va J.K. Rowling tro thanh ti phu.\\
\textit{(The books sold over 500 million copies worldwide, translated into 80 languages, and J.K. Rowling became a billionaire.)}

\end{truyen}
\begin{baihoc}
  Tu choi khong phai la hanh trinh ket thuc --- no la mot trong nhung phan can thiet cua trieu ly thanh cong.\\
  \textit{(Rejection is not the end of the journey --- it is one of the necessary checkpoints on the road to success.)}
\end{baihoc}

\section{Abraham Lincoln --- Dong Thoi Gian Lon Nhat Ve That Bai}
\begin{truyen}{Abraham Lincoln --- Dong Thoi Gian Lon Nhat Ve That Bai}{Abraham Lincoln --- Hoa Ky 1860}
\chuhoa{N}{}N am 1816, gia dinh Lincoln bi mat nha. Nam 1818, me ong mat. Nam 1831, ong kinh doanh phat san.\\
\textit{(In 1816 Lincoln's family lost their home. In 1818 his mother died. In 1831 he failed in business.)}

Nam 1832, ong that cu chuc nhat dinh vien. Nam 1833, kinh doanh phat san lan hai. Nam 1835, hon the chet.\\
\textit{(In 1832 he was defeated for the state legislature. In 1833 his business failed again. In 1835 his fiancée died.)}

Nam 1836, sup do than kinh. Nam 1838, that cu chuc chinh. Nam 1843, that cu vao Quoc Hoi. Nam 1848, mat ghe Quoc Hoi.\\
\textit{(In 1836 he had a nervous breakdown. In 1838 defeated for speaker. In 1843 defeated for Congress. In 1848 lost re-election.)}

Nam 1849, bi tu choi bo nhiem. Nam 1854, that cu vao Thu'ong Vien. Nam 1856, that cu Pho Tong Thong.\\
\textit{(In 1849 he was rejected for a land officer position. In 1854 defeated for Senate. In 1856 defeated for Vice President.)}

Nam 1858, one lan nua that cu vao Thuong Vien --- nam 1860, ong tro thanh Tong Thong My.\\
\textit{(In 1858 defeated for Senate once more --- in 1860 he became President of the United States.)}

\end{truyen}
\begin{baihoc}
  That bai khong tich luy thanh nguoi that bai --- chi tich luy thanh kinh nghiem, neu ban khong buong tay.\\
  \textit{(Failures do not accumulate into a failed person --- they accumulate only into experience, if you do not let go.)}
\end{baihoc}

\section{Nick Vujicic --- Khong Tay Khong Chan Nhung Co Hi Vong}
\begin{truyen}{Nick Vujicic --- Khong Tay Khong Chan Nhung Co Hi Vong}{Nick Vujicic --- Australia 1982}
\chuhoa{N}{}N ick Vujicic sinh ra khong co tay va chan --- mot benh tam co roi rac chua biet nguyen nhan.\\
\textit{(Nick Vujicic was born without arms and legs --- a rare disorder of unknown cause.)}

O tuoi len 8, ong tu hoi tai sao ong phai song --- o tuoi 10, ong da co ke hoach tu tu.\\
\textit{(At age 8 he asked why he should live --- at age 10 he had planned to end his life.)}

Me ong cho ong xem bai bao ve mot nguoi khuyet tat song day du y nghia --- lan dau tien ong thay tu lai minh trong nguoi khac.\\
\textit{(His mother showed him a newspaper article about a disabled person living a full life --- for the first time he saw his reflection in someone else.)}

Ong hoc boi, hoc danh golf, hoc gianh bong, hoc lap trinh may tinh --- tat ca chi bang hai ngon chan nho.\\
\textit{(He learned to swim, play golf, catch a ball, use a computer --- all using only his two small toes.)}

Ong tro thanh dien gia truyen cam hung duoc moi den hon 50 quoc gia, la chong va la cha cua bon con.\\
\textit{(He became a motivational speaker invited to over 50 countries, a husband and father of four children.)}

\end{truyen}
\begin{baihoc}
  Dieu gi ban khong co --- neu dat truoc sau mat minh --- co the tro thanh nguon suc manh lon nhat rather than ra nhung gi ban co.\\
  \textit{(What you lack --- if placed before you as purpose --- can become the greatest source of strength, even greater than what you have.)}
\end{baihoc}

\section{Stephen Hawking --- Vat Ly Hoc Tu Xe Lan}
\begin{truyen}{Stephen Hawking --- Vat Ly Hoc Tu Xe Lan}{Stephen Hawking --- Anh 1963}
\chuhoa{O}{}O tuoi 21, Stephen Hawking bien dang xu phong than kinh ben --- bac si du doan ong chi song them 2 den 3 nam.\\
\textit{(At age 21, Stephen Hawking was diagnosed with ALS --- doctors predicted he would live only 2 to 3 more years.)}

Ong khong chet sau 3 nam --- ong song them 55 nam nua, viet 15 quyen sach va cach mang trong vat ly hoc.\\
\textit{(He did not die after 3 years --- he lived 55 more years, wrote 15 books, and revolutionized physics.)}

Ong mat dan kha nang van dong, roi kha nang noi --- nhung khong mat kha nang suy nghi.\\
\textit{(He gradually lost his ability to move, then to speak --- but never lost the ability to think.)}

Cuoi cung ong giao tiep chi qua cai dung ma ong co the kiem soat --- mot co bang mat --- voi phan mem doc suy nghi.\\
\textit{(In the end he communicated only through the single muscle he could control --- a cheek muscle --- with software reading his thoughts.)}

Ong noi: 'Dieu duy nhat toi su'ng sot hon ca cai mo'i bán' --- su kien ngac nhien la minh van so'ng va lam viec duoc.'\\
\textit{(He said: 'The one thing that surprised me more than my diagnosis --- was the surprise that I could still live and work.')}

\end{truyen}
\begin{baihoc}
  Chang co dieu kien nao tu dong cua het cuoc song --- chi la ban tu dong cua no; hay de tat ca cac cua khac con mo.\\
  \textit{(No condition automatically closes off life --- only you can do that; keep all the other doors open.)}
\end{baihoc}

\section{Malala Yousafzai --- Tien Vao Truong Hoc Bat Chap Dan}
\begin{truyen}{Malala Yousafzai --- Tien Vao Truong Hoc Bat Chap Dan}{Malala Yousafzai --- Pakistan 2012}
\chuhoa{N}{}N am 2012, khi Malala Yousafzai 15 tuoi, mot ten ban vao xe buyt hoc sinh va ban thang vao dau ba.\\
\textit{(In 2012, when Malala Yousafzai was 15, a gunman boarded her school bus and shot her directly in the head.)}

Ba song sot --- sau ca thang phau thuat phuc hoi tai Pakistan va Anh.\\
\textit{(She survived --- after months of surgery and rehabilitation in Pakistan and England.)}

Khi tinh day trong benh vien, cau hoi dau tien cua ba la: 'Chung ban toi vi gi?' Nguoi cha tra loi: 'Vi di hoc.'\\
\textit{(When she awoke in the hospital, her first question was: 'Why did they shoot me?' Her father answered: 'For going to school.')}

Ba tra loi: 'Vay thi toi se tiep tuc di hoc --- vi chinh dieu ho so hai moi la thu quan trong nhat.'\\
\textit{(She replied: 'Then I will keep going to school --- because the very thing they fear is what matters most.')}

Nam 2014, o tuoi 17, ba nhan Giai Nobel Hoa Binh --- nguoi tre nhat trong lich su duoc ghi nhan.\\
\textit{(In 2014, at age 17, she received the Nobel Peace Prize --- the youngest ever to be recognized.)}

\end{truyen}
\begin{baihoc}
  Doi khi chinh vien dan nhung ke co quyen tru ban la luat chung minh rang viec ban dang lam la thu gi do rat quan trong.\\
  \textit{(Sometimes the very bullets those in power aim at you are proof that what you are doing truly matters.)}
\end{baihoc}

\section{Thomas Edison --- 10.000 Thu Nghiem Truoc Den Dien}
\begin{truyen}{Thomas Edison --- 10.000 Thu Nghiem Truoc Den Dien}{Thomas Edison --- Hoa Ky 1878}
\chuhoa{T}{}T homas Edison khong tot nghiep trung hoc --- thay giao noi ong 'qua do' de hoc.\\
\textit{(Thomas Edison never finished high school --- his teacher said he was 'too stupid' to learn.)}

Ong mat viec lam hai lan trong thap nien dau --- nhung moi lan bi sa thai, ong chi nghi: 'Bay gio toi co them thoi gian de thuc nghiem.'\\
\textit{(He lost two jobs in his early years --- but every time he was fired, he thought only: 'Now I have more time to experiment.')}

O tuoi 97, nguoi ta hoi Edison co khi nao tiếc dieu gi khong --- ong tra loi khong bao gio.\\
\textit{(At 97, someone asked Edison if he ever had regrets --- he said never.)}

'Co ai do tinh cho toi la toi da that bai 10.000 lan. Toi khong bao gio 'that bai' 10.000 lan --- toi hoc 10.000 bai.'\\
\textit{('Someone calculated that I had failed 10,000 times. I never 'failed' 10,000 times --- I learned 10,000 lessons.')}

Cac phat minh cua Edison --- den dien, may hat, may dien hinh --- tao nen nen tang cua nen van minh hien dai.\\
\textit{(Edison's inventions --- the light bulb, the phonograph, the movie camera --- formed the foundation of modern civilization.)}

\end{truyen}
\begin{baihoc}
  Chinh nghich canh la thu luyen kien nhan --- va chinh kien nhan la thu thuoc duy nhat chua duoc su tu bo.\\
  \textit{(Adversity itself is the training of patience --- and patience is the only remedy for the disease of giving up.)}
\end{baihoc}

\section{Nguyen Ngoc Ky --- Viet Bang Chan}
\begin{truyen}{Nguyen Ngoc Ky --- Viet Bang Chan}{Nguyen Ngoc Ky --- Viet Nam 1956}
\chuhoa{N}{}N guyen Ngoc Ky bi liet co tay tu nho --- luc dau tuong rang cuoc doi ong se la mot bi kich.\\
\textit{(Nguyen Ngoc Ky was paralyzed in both arms from childhood --- at first it seemed his life would be a tragedy.)}

Nhung ong quyet dinh di hoc --- va hoc bang each dung chan viet chu.\\
\textit{(But he decided to go to school --- and to write using his feet.)}

Ban dau chu ong ngoay ngoac khong doc duoc; nhung qua hang gio luyen viet moi ngay, dan dan chu cai hinh thanh ro rang.\\
\textit{(At first his letters were illegible; but through hours of daily practice, gradually clear characters formed.)}

Ong khong chi hoc xong pho thong ma con vao dai hoc su pham, tro thanh giao vien nhan dang cua mot the he.\\
\textit{(He not only finished school but attended teacher's college, becoming an icon of a generation.)}

Ong noi: 'Toi khong co doi ban tay --- nhung toi co doi ban chan va mot cai dau. Vay la du de song va cong hien.'\\
\textit{(He said: 'I don't have two hands --- but I have two feet and a mind. That is enough to live and contribute.')}

\end{truyen}
\begin{baihoc}
  Ban khong can day du de tro nen phi thuong --- ban chi can su dung het nhung gi ban co.\\
  \textit{(You don't need to be complete to become extraordinary --- you only need to fully use what you have.)}
\end{baihoc}

\section{Florence Nightingale Chong Lai He Thong Y Te}
\begin{truyen}{Florence Nightingale Chong Lai He Thong Y Te}{Florence Nightingale --- Anh 1854}
\chuhoa{T}{}T rong chien tranh Crimea, Florence Nightingale den benh vien quan su tai Scutari va bi sock boi su ban thiu va vo to chuc.\\
\textit{(During the Crimean War, Florence Nightingale arrived at the military hospital in Scutari and was shocked by filth and disorganization.)}

Cac bac si quan su chong doi ba --- mot phu nu khong the noi chuyen ve y khoa trong quan doi.\\
\textit{(Military doctors opposed her --- a woman could not speak about medicine in the army.)}

Ba khong tranh luan bang loi noi --- ba lam viec bang so lieu: thu thap du lieu, ve bieu do, chung minh rang su ban thiu giet chet nguoi.\\
\textit{(She didn't argue with words --- she worked with numbers: collecting data, drawing charts, proving that filth was killing people.)}

Voi du lieu trong tay, ba thuyet phuc Quoc Hoi Anh cai to toan bo he thong y te quan su.\\
\textit{(With data in hand she persuaded the British Parliament to reform the entire military medical system.)}

Ba khong thua truoc nghich canh --- ba bien nghich canh thanh ti le phan tram va bieu do de thay doi the gioi.\\
\textit{(She did not yield to adversity --- she transformed adversity into percentages and charts to change the world.)}

\end{truyen}
\begin{baihoc}
  Doi khi nghich canh khong can duoc vuot qua bang y chi --- ma bang tri tue va du lieu duoc su dung dung cach.\\
  \textit{(Sometimes adversity need not be overcome by willpower --- but by intelligence and data used skillfully.)}
\end{baihoc}
